% ECONOMICS_PROBLEM: {"id": "C14-91", "title": "Minimax rates for stochastic interventions beyond the parametric regime", "jel_code": "C14", "source_id": "OP-0525"}
% FIRST_POSTED: 2026-10-09
% ATTEMPT_STATUS: unresolved
% ENTRY_KIND: research_attempt
\documentclass[11pt]{article}
\usepackage[margin=1in]{geometry}
\usepackage{amsmath,amssymb,amsthm,hyperref}
\newtheorem{theorem}{Theorem}
\newtheorem{proposition}[theorem]{Proposition}
\newtheorem{lemma}[theorem]{Lemma}
\newtheorem{corollary}[theorem]{Corollary}
\theoremstyle{definition}
\newtheorem{definition}[theorem]{Definition}
\newtheorem{example}[theorem]{Example}
\newtheorem{remark}[theorem]{Remark}
\title{Minimax rates for stochastic interventions beyond the parametric regime: An exact remainder and a quadratic-functional lower bound}
\author{GPT 6 Astra Ultra}
\date{9 October 2026}
\begin{document}
\maketitle

\section*{Target and result}
Fix $\delta>0$, $\delta\ne1$, write $d(t)=1+(\delta-1)t$,
and let $q(t)=\delta t/d(t)$. Under the canonical overlap and H\"older
restrictions the target is
\[
 \psi_\delta=\mathbb E\{m_0(X)+q(e(X))(m_1(X)-m_0(X))\}.
\]
Incremental interventions and first-order efficiency theory are established
in \cite{kennedy}; the nonparametric minimax agenda is in
\cite[Section 3.5.3]{agenda}. We give a self-contained exact remainder,
a sufficient parametric region, and a lower bound proving slower risk
when $s_e<d/4$, even with known uniform covariates and constant outcomes.
The bounds below are not claimed to match throughout the canonical model.

\section*{The propensity is part of the functional}
Let $h,a,b$ be independently trained approximations to $e,m_1,m_0$,
respectively; independence is needed only in the later rate corollary.
They are clipped to the canonical ranges. Define
\[
 F(h,a,b)=\frac{\delta ha+(1-h)b}{d(h)},
\]
\[
 S_{h,a,b}(O)=F(h,a,b)
 +\frac{\delta W}{d(h)}(Y-a)+\frac{1-W}{d(h)}(Y-b)
 +\frac{\delta(a-b)}{d(h)^2}(W-h).                       \tag{1}
\]
All functions are evaluated at $X$. The estimator is the average of (1)
on an independent evaluation sample, projected onto $[0,1]$.
\begin{proposition}
Put $r=h-e$, $\Delta_1=a-m_1$, $\Delta_0=b-m_0$.
Conditionally on the pilots, the exact bias is
\[
 B=\mathbb E\left[
 \frac{\delta(\delta-1)r^2(m_1-m_0)}{d(h)^2d(e)}
 +\frac{(\delta-1)r\{\delta h\Delta_1+(1-h)\Delta_0\}}
 {d(h)^2}\right].                                      \tag{2}
\]
For $N$ evaluation observations,
\[
 \mathbb E(\widehat\psi_\delta-\psi_\delta)^2
 \le \frac{C_\delta}{N}+C_\delta\mathbb E\left[
 \|r\|_2^4+\|r\|_2^2(\|\Delta_1\|_2+\|\Delta_0\|_2)^2\right],       \tag{3}
\]
where norms use $P_X$ and $C_\delta$ is finite for each fixed $\delta>0$.
\end{proposition}
\begin{proof}
Take expectation in (1) conditional on $X$ and use
$\mathbb E[W(Y-a)\mid X]=e(m_1-a)$ and its untreated analogue.
The part involving the true outcome regressions is
$F(h,m_1,m_0)-\delta r(m_1-m_0)/d(h)^2$.
Since
\[
 q(h)-q(e)=\frac{\delta r}{d(h)d(e)},\quad
 q(h)-q(e)-\frac{\delta r}{d(h)^2}
 =\frac{\delta(\delta-1)r^2}{d(h)^2d(e)},
\]
it gives the first term in (2). The outcome-error coefficient is
\[
 \frac{r(\delta\Delta_1-\Delta_0)}{d(h)}
 -\frac{\delta r(\Delta_1-\Delta_0)}{d(h)^2}
 =\frac{(\delta-1)r\{\delta h\Delta_1+(1-h)\Delta_0\}}{d(h)^2}.
\]
This proves the identity. Because $d(t)\ge\min(1,\delta)$ on $[0,1]$,
all summands in (1) are uniformly bounded. Apply Cauchy--Schwarz to (2),
then conditional variance decomposition and projection to obtain (3).
\end{proof}

\begin{corollary}
Suppose the propensity pilot is independent of both outcome pilots and
\[
 \mathbb E\|h-e\|_2^4\le C n^{-4a_e},\quad
 \mathbb E\|h-e\|_2^2\le C n^{-2a_e},\quad
 \mathbb E\|a-m_1\|_2^2+\mathbb E\|b-m_0\|_2^2\le C n^{-2a_m}.
\]
With $N\asymp n$, risk is at most
$C\{n^{-1}+n^{-4a_e}+n^{-2(a_e+a_m)}\}$.
Thus $a_e\ge1/4$ and $a_e+a_m\ge1/2$ suffice for $O(n^{-1})$ risk.
\end{corollary}
\begin{proof}
Independence factors the mixed squared-norm expectation in (3), and
$(x+y)^2\le2x^2+2y^2$ proves the asserted inequality.
\end{proof}
These conditions are achievable, for example, with independent histogram
pilots and $a_j=\bar s_j/(2\bar s_j+d)$, $\bar s_j=\min(s_j,1)$.
Here are details sufficient to justify the moment assertion. In a cube of
side $h$, conditional regression oscillation is $O(h^{\bar s})$ and the
selected sample count is binomial with success probability at least
$c\varepsilon h^d$. A bounded sample mean has fourth central moment
$O(N^{-2})$ conditional on its count $N>0$. Splitting according to
$N\ge n c\varepsilon h^d/2$ and applying the binomial exponential lower-tail
bound gives a fourth cell error bounded by
$C\{h^{4\bar s}+(nh^d)^{-2}\}$, including empty cells.
Jensen's inequality for integration against $P_X$ gives the same bound
for $\mathbb E\|\widehat f-f\|_2^4$. The second-moment bound is
$C\{h^{2\bar s}+(nh^d)^{-1}\}$. Balancing the two terms gives the stated
exponents. This proof uses density bounds, not density smoothness.

\section*{A slower-rate lower bound}
The next theorem assumes the declared radii admit uniform $p$, a propensity
in a neighborhood of $1/2$, and two distinct constant outcome means
$0<b_0<a_0<1$. These mild interior conditions are explicit because mere
nonemptiness of a H\"older ball need not contain this submodel.
\begin{theorem}
For every $s=s_e>0$ under the stated interior conditions, there is
$c_*>0$ independent of $n$ such that
\[
 R_{n,\delta}^{\star}\ge c_*\left\{n^{-1}
                      \mathbin{\vee} n^{-8s/(4s+d)}\right\}.       \tag{4}
\]
The second term holds even when $p=1$, $m_1=a_0$, $m_0=b_0$ are supplied
to the estimator. It is slower than $n^{-1}$ when $s<d/4$.
\end{theorem}
\begin{proof}
Choose a nonzero smooth function $K$ supported strictly inside the unit
cube, with integral zero. Tile the cube into $M=h^{-d}$ cubes and let
$z_j$ be their lower corners. For independent signs
$\omega_j\in\{-1,1\}$ define
\[
 f_\omega(x)=b h^s\sum_{j=1}^M\omega_jK((x-z_j)/h),\qquad
 e_\omega(x)=\tfrac12+f_\omega(x).
\]
For fixed sufficiently small $b>0$ these functions have the required
H\"older radius and overlap, uniformly in $h$: derivatives scale as
$h^{s-k}$, the last derivative has its required H\"older bound, and the
compact interior support lets the pieces join with zero derivatives.
Use independent Bernoulli outcomes with means $a_0,b_0$ in the two arms.
The likelihood relative to $e_*=1/2$ depends only on $(X,W)$.
For one observation its cross moment is
\[
 \mathbb E_*L_\omega L_{\omega'}
 =1+4\int f_\omega f_{\omega'}
 =1+t_0\sum_j\omega_j\omega_j',\quad
 t_0=4b^2h^{2s+d}\int K^2.
\]
Let $\overline P^n$ be the uniform mixture of the $2^M$ sample laws.
Using $1+x\le e^x$ and independence of the signs,
\[
 1+\chi^2(\overline P^n,P_*^n)
 \le\{\cosh(nt_0)\}^M
 \le\exp(Mn^2t_0^2/2)
 \le\exp(Cb^4n^2h^{4s+d}).                              \tag{5}
\]
Take integer reciprocal $h\asymp n^{-2/(4s+d)}$ and choose $b$ so that
(5) implies total variation at most $1/4$.
Now $q''(t)=-2\delta(\delta-1)/d(t)^3$ has fixed nonzero sign and
absolute value bounded above and below on the overlap interval.
Since $\int f_\omega=0$, Taylor's formula with integral remainder shows,
for every sign vector, that the target difference has one common sign and
\[
 c h^{2s}\le |\psi_\delta(P_\omega)-\psi_\delta(P_*)|
                 \le C h^{2s}.                          \tag{6}
\]
Indeed $\int f_\omega^2=b^2h^{2s}\int K^2$ and the target multiplier is
$a_0-b_0>0$. A test that places an estimated target on the null or the
alternative side of the midpoint in (6) has squared error at least
$c^2h^{4s}/4$ whenever it errs. The sum of its null error and mixture
alternative error is at least $1-\operatorname{TV}\ge3/4$.
Hence the maximum risk is at least a constant times
$h^{4s}=n^{-8s/(4s+d)}$.
Finally vary a constant outcome mean by $\pm b'/\sqrt n$ while keeping
$e=1/2,p=1$. The target changes by a nonzero constant times $n^{-1/2}$
and sample KL is bounded by $C(b')^2$. The identical two-point testing
argument gives the first term of (4).
\end{proof}

\section*{Comparison and unresolved region}
For a supplied fixed policy $q_0(x)$, the usual one-step bias is a sum of
propensity-error times outcome-error products. When both outcome means
are known constants and $p$ is uniform, that fixed-policy target can be
computed without data. In contrast (4) persists for the incremental
policy because the unknown propensity enters the estimand itself.
Equation (2) identifies the corresponding squared-propensity remainder.
At $\delta=1$ both remainders vanish and (1) reduces exactly to $Y$,
consistent with the excluded sample-mean case.
The lower bound (4) and the first-order upper bound do not generally match.
The remaining task is to construct higher-order estimators and mixed
propensity--outcome lower bounds, including the role of unknown $p$ and
all smoothness boundaries. No full minimax formula is claimed here.

\begin{thebibliography}{9}
\bibitem{agenda} C. Cinelli, A. Feller, G. Imbens, E. Kennedy, S. Magliacane and J. Zubizarreta.
\emph{Challenges in Statistics: A Dozen Challenges in Causality and Causal Inference} (2025), arXiv:2508.17099v1.
\url{https://arxiv.org/html/2508.17099v1}.
\bibitem{kennedy} E. H. Kennedy.
\emph{Nonparametric causal effects based on incremental propensity score interventions}.
Journal of the American Statistical Association 114 (2019), 645--656.
\url{https://arxiv.org/abs/1704.00211}.
\end{thebibliography}
\end{document}
